%% notes-tex/030-scaling-plots/030-scaling-plots.tex
%% [[experiments.030-solid-growth-multi.scripts.scaling_plot_catalog_030]]
%% https://github.com/Mjvolk3/torchcell/tree/main/notes-tex/030-scaling-plots/030-scaling-plots.tex
%%
%% Build:  make          -> 030-scaling-plots.pdf
%%         make check    -> formatting / provenance / citation gate
%%         make figures  -> re-export the four draw.io sketches (DRAWIO=<binary>)
%%
%% Table 1: experiments/030-solid-growth-multi/scripts/scaling_plot_catalog_030.py.
%% Figures: notes/assets/drawio/030-scaling-plots.gen.py (sketches, no data).
%% Nothing in this document is a measurement; every curve is a shape the plot would have.

\documentclass[10pt,a4paper]{article}
\usepackage[draft]{tcdoc}
\usepackage{float}
\usepackage{amsmath}
\usepackage{amssymb}

\emergencystretch=3.5em

\title{\vspace{-1.5em}Plots the 030 Scaling Grid Can Feed\\
  {\large Sixteen plots, the runs and keys each one draws on, and what each
   would show, sketched before any run has reported}}
\author{Michael Volk}
\date{\today}

\begin{document}
\maketitle
\thispagestyle{empty}

\begin{abstract}
\noindent
The forty jobs of the 030 scaling grid produce, per run and per epoch, a
validation loss in label units pooled and per screen, the records seen, the
measured FLOPs and the GPU hours, and per point two reads of that loss: the
validation-selected epoch and an interpolation at one shared records-seen budget.
Those rows feed sixteen plots in four families. The surface family reads the
loss across the five-by-four grid both ways, against data at every size and
against size at every subset, as a heatmap, and as the fitted forms on $L - E$.
The compute family draws the learning curves against records seen, FLOPs and GPU
hours, whose lower envelope is the compute-efficient frontier, and the cost per
record and pace against size from the calibration jobs. The diagnostic family is
what says how far an exponent can be trusted: the bootstrap distribution, the
held-out prediction against the observation, the fitted floor against the
replicate floor, the seed spread against the step between cells, the
learning-rate check, and the per-screen Pearson beside the loss. The pair family
is the 030 minus 025 gap and the paired learning curves. Nine plots the queued
grid feeds in full and seven in part, the part being a stated limit rather than a
missing row: ten-epoch full-pool points, one seed at most cells, four points on
the data axis. One catalog table names every plot, its axes, its read, the rows it
needs and its status, and four sketches show the shape each family would take.
\end{abstract}

\vspace{0.5em}
\tccontents
\clearpage

%% notes-tex/030-scaling-plots/sections/1-inputs.tex

\section{What the runs provide\secstatus{todo}}
\label{sec:inputs}

\begin{keybox}
\textbf{Measured.} The rows each run logs and the reads the readout takes from
them are implemented and tested (\texttt{scaling\_readout\_030.py},
\texttt{torchcell.trainers.compute\_accounting}). The grid is queued on Delta.

\textbf{Not measured.} Every curve in every sketch. No run has reported, and no
plot here has a data point behind it.
\end{keybox}

The grid, its reads and its limits are set out in the companion study
(\texttt{notes-tex/030-scaling-study}), whose vocabulary applies here. In short:
the \term{loss} is the validation mean squared error on the gene-interaction
score in \term{label units}, after the inverse of the training normalization,
pooled over the pinned 37,705 triples and per \term{screen} (Kuzmin 2018,
Kuzmin 2020) through the per-token keys; $N$ is non-embedding parameters; $D$
is distinct training records, and on 030 also gene-interaction entry rows; the
\term{selected read} is the loss at the epoch of lowest pooled validation MSE;
the \term{matched read} is the loss interpolated at one shared records-seen
budget $B$, the smallest total any point of the grid reached; the
\term{replicate floor} $E = (1 - r)\,\mathrm{var}(\tau)$ is the loss no predictor
of a single measurement can beat, per screen; the \term{held-out test} refits
without the largest point of an axis and predicts it.

\notesec{Rows}
Each run logs, at the end of every epoch, the pooled and per-screen validation
MSE and Pearson, the records seen, the cumulative measured FLOPs and GPU hours,
and the epoch index that joins them. A point is one (size, subset, seed, learning
rate); a chain of jobs is one point with its epochs stitched by the checkpoint
each link resumed from. The readout keeps every epoch of every point
(\texttt{<sweep>\_epochs.csv}), one row per point and read
(\texttt{<sweep>\_points.csv}), every fit with its interval and held-out
prediction (\texttt{<sweep>\_fits.json}), and the gap per subset
(\texttt{<sweep>\_gap.csv}). Every plot below is drawn from one of those four
files, so a plot is regenerated from the rows and never assembled by hand.

\notesec{The catalog}
Table~\ref{tab:catalog} is the list. The id is the panel letter family used in
the sketches: S for the loss surface, C for compute, D for diagnostics, P for the
pair. Status is read against the queued grid as it stands: full means every row
the plot needs exists at the design's seed count; partial means the rows exist
but a limit of the design narrows what the plot can claim, named in its section.

%% GENERATED FILE -- do not hand-edit.
%% SOURCE: experiments/030-solid-growth-multi/scripts/scaling_plot_catalog_030.py (the PlotSpec catalog)
\begin{table}[htbp]
\centering\footnotesize
\caption{The plot catalog: 16 plots the queued grid can feed, 9 in full and 7 in part. Read names what the y value is at each point: per-epoch rows, the validation-selected epoch, the interpolation at the shared records-seen budget $B$, or the calibration jobs. Feeds names the rows of the grid the plot needs. Partial means the rows exist but a stated limit narrows the claim (ten-epoch full-pool points, one seed, four points on an axis).}
\label{tab:catalog}
\begin{tabular}{llp{0.30\linewidth}p{0.16\linewidth}p{0.26\linewidth}l}
\toprule
id & family & plot (x; y) & read & feeds & status \\
\midrule
S1 & surface & \textbf{Loss against D at every size}: D, training records (log); 030 also in entry rows; val MSE, label units (log) & both & the 5 x 4 seed-1 grid; seeds 2-3 at the two smallest subsets for 45x2 and 180x8 & full \\
S2 & surface & \textbf{Loss against N at every subset}: N, non-embedding parameters (log); val MSE, label units (log) & both & the 5 x 4 seed-1 grid & partial \\
S3 & surface & \textbf{The joint surface as a 5 x 4 heatmap}: D, four subsets; N, five sizes & both & the 5 x 4 seed-1 grid & full \\
S4 & surface & \textbf{Fitted forms on the two axes}: D or N (log); val MSE minus the fitted E (log) & both & S1's 180x8 column and S2's full-pool column & partial \\
C1 & compute & \textbf{Learning curves against records seen}: records seen (log); val MSE, label units (log) & per-epoch & every run; chains stitched by resumed\_from & full \\
C2 & compute & \textbf{Loss against training FLOPs, the envelope}: cumulative FLOPs (log); val MSE, label units (log) & per-epoch & every run on the full pool and on 1/4 (the large-data columns) & partial \\
C3 & compute & \textbf{Loss against GPU hours}: GPU hours (log); val MSE, label units (log) & per-epoch & every run & full \\
C4 & compute & \textbf{Cost per record and pace against size}: N (log); FLOPs per record (left); records per second (right) & calibration & the two calibration jobs (180x8, 45x2) now; every grid run's first epoch after & partial \\
D1 & diagnostic & \textbf{Bootstrap distribution of each exponent}: exponent value; count of resamples & both & 300 resamples over runs of S1's and S2's fitted columns & full \\
D2 & diagnostic & \textbf{Held-out prediction against observation}: predicted loss at the held-out point; observed loss & both & every fit of S4 & partial \\
D3 & diagnostic & \textbf{Fitted E against the replicate floor}: screen (2018, 2020, pooled); loss, label units & both & S4's floored fits and the noise-ceiling file & full \\
D4 & diagnostic & \textbf{Seed spread against effect size}: (size, subset) cell; val MSE, label units & selected & the eight seed-replicate cells plus their seed-1 neighbors & full \\
D5 & diagnostic & \textbf{The learning-rate check}: epoch; val MSE, label units (log) & per-epoch & the two lr jobs and the seed-1 cell at 45x2 on 1/4 & partial \\
D6 & diagnostic & \textbf{Per-screen Pearson beside the loss}: D (log); val Pearson, per screen and pooled & selected & the 180x8 column & full \\
P1 & pair & \textbf{The 030 minus 025 gap against D}: D, training records (log); 030 minus 025 pooled val MSE, label units & both & the 180x8 column of both sweeps & partial \\
P2 & pair & \textbf{Paired learning curves}: records seen (log); val MSE, label units (log) & per-epoch & the 180x8 column of both sweeps & full \\
\bottomrule
\end{tabular}
\end{table}


%% notes-tex/030-scaling-plots/sections/2-surface.tex

\section{The loss surface, read both ways\secstatus{todo}}
\label{sec:surface}

The grid is a five-by-four sample of $L(N, D)$ at one seed, with the two end
sizes replicated at the two smallest subsets. That is what makes the surface
family possible: every row of Figure~\ref{fig:surface} is a cut of the same
surface, and the two single-axis fits of the study are its edge columns.

\tcfig{figures/030-scaling-plots-surface.pdf}{%
The surface family, sketched. \textbf{a}, S1: loss against $D$ at every size,
one curve per size. Parallel curves on log axes are one data exponent for the
family; a fan is a data exponent that depends on size, which the single-column
fit of the study would not see. The replicate floor is dashed. \textbf{b}, S2:
loss against $N$ at every subset, one curve per subset. Each curve flattens at
the plateau its data sets, $E + B D^{-\beta}$, dashed per curve, and the size at
which a curve meets its plateau is the size beyond which more parameters stop
paying at that $D$. \textbf{c}, S3: the same twenty cells as a heatmap of the
matched read, with a second heatmap for the selected read in the real figure.
\textbf{d}, S4: the fitted forms drawn on $L - E$, where the floored form is a
straight line, with the bootstrap band and the held-out point ringed. Every
curve is hand-drawn to show the shape the plot would hold.
}{fig:surface}
\src{notes/assets/drawio/030-scaling-plots.gen.py}

\notesec{S1, loss against data at every size}
Five curves, each a row of the grid, on both reads and per screen. The study fits
the 180x8 curve alone; the other four are the control on that fit. If the small
sizes' curves run parallel to it, the data exponent is a property of the data
and the single fit stands for the family. If the small curves flatten earlier,
the small sizes are capacity-limited at large $D$, and the exponent fit at 180x8
is the largest the family shows, not the family's. The 1/64 and 1/16 points of
the end sizes carry three seeds each, so the spread of the replicated cells is
drawn as the noise floor of every curve.

\notesec{S2, loss against size at every subset}
Four curves, each a column. On the matched read the curves are directly
comparable, since every cell is read at the same records seen. On the selected
read the full-pool curve is a ten-epoch slice and sits above where convergence
would put it, so the plateau it appears to meet is partly the budget, which is
the limit that keeps this plot partial. The plateau lines are the second term of
the floored form fitted on each column, and the gap between a plateau and the
replicate floor is what more data could still buy at that size.

\notesec{S3, the heatmap}
Twenty cells shaded by the matched read, with the selected read beside it. The
reading is local: at any cell the steeper descent, down the column toward
larger $N$ or along the row toward larger $D$, is the direction the next unit of
effort should take from that cell. The two heatmaps differ most at the full-pool
column, and the size of that difference is a picture of how far the ten-epoch
budget sits from convergence.

\notesec{S4, the fitted forms}
Each axis of the study drawn on $L - E$ for the floored form, so the fit is a
straight line, with the pure power law beside it and the 5 to 95 percent band
from 300 resamples over runs. The held-out largest point is ringed with the
prediction the in-sample fit made for it. On the $D$ axis there are four points,
so the floored held-out test is undefined (three points fit three parameters)
and only the power form is tested there; that is the partial status.

%% notes-tex/030-scaling-plots/sections/3-compute.tex

\section{Compute\secstatus{todo}}
\label{sec:compute}

The compute callback is what lets the grid say anything about $C$. It measures
FLOPs on five real training steps of every run, accumulates them per epoch with
the records seen and the GPU hours, and saves its state with the checkpoint, so
a chained point's totals continue. Figure~\ref{fig:compute} is what those keys
draw.

\tcfig{figures/030-scaling-plots-compute.pdf}{%
The compute family, sketched. \textbf{a}, C1: per-epoch validation loss against
records seen, one line per point colored by size, one panel per subset in the
real figure; the dot on each line is its selected epoch and the vertical line is
the matched budget $B$. \textbf{b}, C2: the same trajectories against cumulative
FLOPs; the dashed lower envelope is the compute-efficient frontier, the size
that reaches a given loss for the least compute. \textbf{c}, C3: the same against
GPU hours, the unit the queue charges; the horizontal offset of a size between b
and c is its utilization. \textbf{d}, C4: FLOPs per record and records per
second against size from the calibration jobs, with the spread over the five
measured steps as a whisker; a flat pace line across sizes means the loader, not
the GPU, sets the pace.
}{fig:compute}
\src{notes/assets/drawio/030-scaling-plots.gen.py}

\notesec{C1, learning curves against records seen}
The plot every other read is taken from. It shows where each point's selected
epoch sits, where $B$ cuts each curve, and which curves are still falling at
their last epoch. A full-pool curve that is still falling at ten epochs is the
picture of the unconverged headline point, and it is the plot to show beside any
exponent fit on the selected read.

\notesec{C2 and C3, the envelope}
Against FLOPs, the lower envelope of the trajectories across sizes is the
compute-efficient frontier of this architecture on this data: at each budget the
size whose curve is lowest. The grid has one budget line of an IsoFLOP sweep,
the fixed-epoch column, so the envelope here is read as the shape of the
frontier rather than its slope, which keeps C2 partial. Against GPU hours the
envelope is the same picture in the unit the queue charges, and because the FLOP
count excludes the loader while the wall clock includes it, the ratio of a size's
position in C2 to its position in C3 is the utilization that size achieved.

\notesec{C4, cost and pace against size}
From the two calibration jobs now, and from every grid run's first epoch after.
FLOPs per record should rise with $N$; whether records per second falls with it
is the question the calibration was queued to answer. A pace that follows the
FLOPs means the GPU binds and the small sizes could have trained longer in their
48 h; a pace that stays flat means the loader binds and no size finishes sooner
than 180x8. The whisker is the minimum to maximum over the five measured steps,
which on the per-entry path is the spread one batch's row count introduces.

%% notes-tex/030-scaling-plots/sections/4-diagnostics.tex

\section{Diagnostics\secstatus{todo}}
\label{sec:diagnostics}

An exponent is a fitted quantity with an interval, and the interval, not the
point, is what an extrapolation inherits. The six plots of
Figure~\ref{fig:diagnostics} are what put the interval and its caveats on the
page beside the curves.

\tcfig{figures/030-scaling-plots-diagnostics.pdf}{%
The diagnostic family, sketched. \textbf{a}, D1: the bootstrap distribution of
one exponent, 300 resamples over runs, the 5 to 95 percent interval shaded.
\textbf{b}, D2: the held-out prediction of every fit against the observed loss
at the held-out point; a marker on the diagonal is a fit that extrapolated one
step. \textbf{c}, D3: the fitted $E$ with its interval against the replicate
floor of each screen, the floor bounds as a band and the mean-prediction loss as
the top of the range. \textbf{d}, D4: the three seeds of each replicated cell
beside the step to its neighbor along $D$. \textbf{e}, D5: the learning-rate
check at 45x2 on 1/4. \textbf{f}, D6: validation Pearson against $D$, pooled
with its 0.35 screen-offset floor and per screen through the per-token keys.
}{fig:diagnostics}
\src{notes/assets/drawio/030-scaling-plots.gen.py}

\notesec{D1, the bootstrap}
One histogram per fitted exponent (axis, read, target, form). On the $D$ axis
with four points and one seed at two of them, the resamples that drop a point
move the slope a long way, and the histogram will be wide; drawing it is the
honest way to say so, and it is the plot that keeps the point estimate from
being quoted alone.

\notesec{D2, held out against observed}
Every fit in one scatter: predicted loss at the held-out largest point against
the observed loss there. The fits that land on the diagonal are the ones whose
exponent extrapolates one step; the ones that miss describe the points they
were fit on and nothing beyond them. The floored form on the $D$ axis has no
marker here, which is why the plot is partial.

\notesec{D3, the fitted floor against the replicate floor}
Per screen and pooled: the fitted $E$ with its interval, the band between the
floor bounds (for 2018, the Dango 0.59 and the cross-year 0.89 reliabilities;
for 2020 the same-strain 0.44), and the loss of predicting the screen mean as the
top of the range. A fitted $E$ inside or just above the band is a model near the
noise limit of that screen; one far above it says the model family, not the
data, is the limit; one below it is a wrong fit. On 2020 the floor is 56\% of
the variance, so the band sits high and the plot will show how little of the
2020 loss any model can remove.

\notesec{D4, seed spread against effect size}
The four replicated cells, three seeds each, with the step to the adjacent cell
along $D$ drawn beside them. Every exponent interval rests on this spread; if a
step along the axis is not larger than the spread within a cell, the slope
between those cells is noise, and the plot says so before a fit is quoted.

\notesec{D5, the learning-rate check}
Three curves at the smallest size on 1/4. It says whether 2.5e-4 is near the
optimum at 45x2 and nothing about 180x8, which runs at a batch four times
smaller than the one it was tuned at; the plot is drawn with that caption so
the $N$ exponent is not read as if the rate had been checked across the ladder.

\notesec{D6, Pearson beside the loss}
The headline metric every reader of the trigenic work expects, drawn with the
floor the two screen means put under the pooled number: predicting only the
per-screen mean scores 0.34 on the validation entries. The per-token Pearson
curves carry no such floor, so this plot and the per-screen loss curves of S1
tell one story, and a pooled Pearson that rises while the per-screen loss does
not fall is the token learning the offset, not the model learning interactions.

%% notes-tex/030-scaling-plots/sections/5-pair.tex

\section{The 025 pair\secstatus{todo}}
\label{sec:pair}

The 025 twin trains the merged-label trainer on the same subsets, the same
genotypes pair for pair, at 180x8, with three seeds at the two smallest
subsets. It passes through the same readout and the same two reads at the same
$B$, so the two plots of Figure~\ref{fig:pair} compare like with like.

\tcfig{figures/030-scaling-plots-pair.pdf}{%
The pair family, sketched. \textbf{a}, P1: the 030 minus 025 pooled validation
MSE per subset, seed mean and sd where both sides have three seeds, a single
point at 1/4, and the full point hollow as the difference of two unconverged
ten-epoch runs. \textbf{b}, P2: the paired learning curves against records seen,
one panel per subset in the real figure; a gap that holds past $B$ is more
supervision per genotype, one that closes is a faster start.
}{fig:pair}
\src{notes/assets/drawio/030-scaling-plots.gen.py}

\notesec{P1, the gap against data}
Pooled only, since 025 logs no per-screen key. Interpretable at 1/64 and 1/16,
where both sides converge and both have a seed spread; a single comparison at
1/4; a learning-speed difference at the full pool. The gap is per-entry rows
plus the dataset token plus the measurement types the 030 build added, and the
caption says so until the token ablation column exists.

\notesec{P2, paired learning curves}
The plot that tells the two readings of a gap apart. If the 030 curve sits below
the 025 curve from the first epochs and the distance holds through $B$ and
beyond, the per-entry rows are more supervision at every stage of training. If
the curves start apart and meet, the extra rows were a faster start that the
merged label catches up on, and the gap at a fixed $B$ overstates what the rows
buy at convergence.

%% notes-tex/030-scaling-plots/sections/6-beyond.tex

\section{What this grid cannot draw\secstatus{todo}}
\label{sec:beyond}

Three plots of a full scaling section are outside what the queued rows provide,
and each is named here so it is not improvised from them.

\notesec{The IsoFLOP curves and the frontier's slope}
A U of loss against $N$ at a fixed compute budget needs several sizes trained
to the same FLOP count, and the frontier needs three such budgets. The grid has
one budget line, the fixed-epoch column, whose envelope in C2 is the frontier's
shape at one budget and nothing about its slope. Three budgets spanning two
decades with four sizes each is the sweep, twelve jobs, after the data and size
sweeps have shown the budget buys anything.

\notesec{Converged loss against size}
S2's full-pool curve is a ten-epoch slice. The converged curve needs the small
sizes run to convergence on the full pool, which the per-size epoch budget after
calibration allows if the GPU binds, and a converged 180x8 point, which is nine
chained 48 h jobs. Until then the $N$ exponent is drawn with its caption and the
heatmap carries both reads so the gap between them is visible.

\notesec{The per-screen gap}
P1 is pooled because the 025 trainer logs no per-screen key. A per-screen gap
needs the 025 validation predictions dumped per epoch and scored against the
screen each triple belongs to, which is a small addition to the 025 trainer and
a rerun of its eight jobs, or a one-off scoring of its checkpoints.


\end{document}
